% !TeX root = ../机器学习笔记.tex
% 由主文档通过 \input 载入；请编译主文档。
\setcounter{chapter}{23} % 保留原书章号。
\chapter{强化学习导论}
\label{chap:ml-rl}

强化学习的行动影响未来状态和可获得的数据，因此决策不能只看当前预测误差。本章采用有限状态动作空间和折扣回报，说明模型、价值、策略和求解算法的关系。

\section{从奖励反馈到序贯决策}

假设两个选项的成功概率未知，试用后只能看到所选项的成功与否。一直选择目前平均奖励较高的选项可能错过尚未充分尝试的更好选项；尝试未知选项可以获得信息，却牺牲即时收益。这是多臂赌博机（multi-armed bandit）中的探索与利用问题。进入迷宫后，动作还改变下一步能到达的位置，问题进一步成为序贯决策。

环境实现生成状态与奖励，数学模型描述这种生成规律，算法利用模型或样本学习，策略规定行动选择。四者不是同一对象。奖励评价已经产生的结果，通常不直接给出当前的正确动作标签；策略变化也会改变后续训练数据分布。

\section{马尔可夫决策过程与策略}

有限马尔可夫决策过程（Markov decision process，MDP）包含状态集 $\mathcal S$、各状态的非空有限动作集、初始分布、转移概率 $P(s'\mid s,a)$、条件期望奖励 $r(s,a)$ 和折扣因子 $0\leq\gamma<1$。假定单步随机奖励有界，且给定当前状态动作后，下一状态奖励的条件分布不再依赖更早历史。状态是满足该建模假设的信息概括，观测则是智能体实际看见的量，部分可观测时两者不能等同。

平稳随机策略 $\pi(a\mid s)$ 对每个状态给出动作概率，确定性策略指定一个动作。回报 $G_t=\sum_{k\geq0}\gamma^kR_{t+k+1}$ 在有界奖励下绝对收敛。终止状态可表示为之后奖励恒为零的吸收状态，以统一有限回合与折扣形式。

\section{价值函数与 Bellman 方程}

价值定义为 $V^\pi(s)=\mathbb E_\pi[G_t\mid S_t=s]$，动作价值为 $Q^\pi(s,a)=\mathbb E_\pi[G_t\mid S_t=s,A_t=a]$。由 $G_t=R_{t+1}+\gamma G_{t+1}$ 和全期望公式，
\begin{equation}
 V^\pi(s)=\sum_a\pi(a\mid s)\left[r(s,a)+\gamma\sum_{s'}P(s'\mid s,a)V^\pi(s')\right].
 \label{eq:ml-bellman-expectation}
\end{equation}
右侧定义 Bellman 算子 $T^\pi$，其在最大范数下满足
$\|T^\pi V-T^\pi W\|_\infty\leq\gamma\|V-W\|_\infty$，
因此固定点唯一。

最优性算子 $TV(s)=\max_a[r(s,a)+\gamma\sum_{s'}P(s'\mid s,a)V(s')]$ 也具有相同压缩性，因为两组数的最大值之差绝对值不超过最大逐项差。其唯一固定点记为 $V^*$。在每个状态选择达到最大值的动作形成策略 $\pi^*$，有 $T^{\pi^*}V^*=TV^*=V^*$，故它的价值就是 $V^*$；由 Bellman 不等式及有界折扣尾项可知其他策略的价值不超过它。

\section{已知模型下的动态规划}

策略评估反复用 $V_{k+1}=T^\pi V_k$，或直接解线性方程 $(I-\gamma P^\pi)V=r^\pi$。策略改进按当前 $V^\pi$ 选择贪心动作；新算子作用于旧价值不低于旧价值，反复迭代并取极限得到新策略价值不低于原策略。精确策略迭代交替评估与改进，并对并列采取保留原动作等固定规则，以避免无意义切换。

价值迭代直接用 $V_{k+1}=TV_k$，压缩性给出 $\|V_k-V^*\|_\infty\leq\gamma^k\|V_0-V^*\|_\infty$。还可由残差得到
$\|V-V^*\|_\infty\leq\|TV-V\|_\infty/(1-\gamma)$。
这些算法需要已知模型；MDP 是问题的数学描述，动态规划是求解它的方法。

\section{从样本学习价值与策略}

不知模型时，可用完整回合回报作 Monte Carlo 估计，也可用时间差分（temporal-difference，TD）目标 $R_{t+1}+\gamma V(S_{t+1})$ 更新
$V(S_t)\leftarrow V(S_t)+\eta_t[R_{t+1}+\gamma V(S_{t+1})-V(S_t)]$。
前者等到后续结果出现，后者使用已有价值估计自举，误差性质不同。

Q-learning 更新
\begin{equation}
 Q(S_t,A_t)\leftarrow Q(S_t,A_t)+\eta_t\left[
 R_{t+1}+\gamma\max_aQ(S_{t+1},a)-Q(S_t,A_t)\right],
 \label{eq:ml-qlearning-update}
\end{equation}
终止转移的自举项为零。行为策略负责采样，目标中的最大值对应贪心控制，故它是异策略方法。这里引用表格情形的收敛结论\cite{ml-qlearning}：在有限固定 MDP、有界奖励、$0\leq\gamma<1$ 下，从有限实数初值出发，若每个状态动作对几乎必然被访问无穷次，步长由更新前的信息确定、取值于 $(0,1)$，且对每对各自的访问序列满足步长和发散、平方和收敛，则估计几乎必然收敛到最优动作价值。该结论不能直接用于神经网络函数逼近。

另一条路线直接优化参数策略。有限时域、固定且与策略参数无关的环境、可微且支持不变的正概率策略下，可由轨迹概率求导得到
$\nabla_\theta\mathbb E_\theta G_0=\mathbb E_\theta[G_0\sum_t\nabla_\theta\log\pi_\theta(A_t\mid S_t)]$，
这里假定求导与有限期望可交换。这说明策略梯度利用回报加权行动概率梯度，而不是学习每一步的最优动作标签。

\section{例题、练习与学习检查}

\begin{example}[即时奖励与长期回报]
起点选动作 A 立即获得二并终止；选动作 B 立即获得零，下一步必定获得五并终止，$\gamma=0.9$。
\end{example}
\begin{solve}
两动作价值分别为二和 $0+0.9\times5=4.5$，故 B 更优，尽管其即时奖励较小。若一次终止转移奖励为二、旧 $Q$ 为零、步长为 $1/2$，Q-learning 更新后该项为一；它是一次样本更新，不是对完整模型的精确求解。
\end{solve}

习题：证明 $T^\pi$ 的压缩性；用两状态模型手算两轮价值迭代；说明离线强化学习仍需处理行动覆盖不足；设计把训练探索回报与固定策略评价回报分开报告的实验。





